%% =====================================================================================
\section{The certified computations}\label{sec:computation}
%% =====================================================================================

\subsection{Arithmetic with directed rounding}\label{sec:rounding}

The certified checkers \path{certificate-14501/src/rigx.c} (for \cref{thm:main}) and \path{src/rigcert.c} (for
\cref{thm:lean}) are single C files that implement \cref{sec:evaluation} in IEEE-754 double precision with directed
rounding. The first was derived from the second and keeps all of its conventions. (Separate implementations in
fixed-point integer arithmetic are part of the formalizations of \cref{thm:main,thm:lean}; see
\cref{sec:certified-formal,sec:lean}.) The conventions
are as follows; \cref{app:rounding} lists the rounding direction of the quantities.
\begin{itemize}
\item The processor is switched to round-upward mode once and stays there, so every $+,-,\times,/$ returns an
  upper bound for the exact result of its operands. Lower bounds are computed as
  $\mathrm{RD}(x\circ y)=-\mathrm{RU}(-(x\circ y))$. Square roots are corrected to the required direction by an
  exact fused-multiply-add sign test.
\item No library transcendental function enters any bound. Powers $q^{-a}$ are formed by repeated division,
  $(a+1)^\theta$ by repeated multiplication and a directed square root, and the logarithms in the test of
  \cref{prop:bbmst61} are replaced by rigorous lower bounds from a truncated $\operatorname{artanh}$ series with
  positive terms. Library functions such as \code{log}, \code{pow} and \code{sqrt} are called only to evaluate the
  schedule and to choose free parameters (enumeration cut-offs, the split point $y$ of \cref{sec:SBH}, table sizes).
\item Quantities that enter a total with a plus sign are rounded up; quantities that are subtracted, such as the
  partial Euler sums in \eqref{eq:EU}, the thresholds $\beta$ in \cref{sec:SBH} and the enumerated mass in
  \eqref{eq:AB-bound}, are rounded down. Differences that would cancel badly are replaced by algebraically identical
  sums of nonnegative terms, such as the lost-mass recursion \eqref{eq:lost} and the suffix sums \eqref{eq:suffix}.
\item The decimal output is formatted from integers, so the printed totals are themselves upper bounds.
\item A start-up self-test aborts if directed rounding or the negation idiom does not behave as expected, or if
  constant folding in round-to-nearest is detected. The sources forbid \code{-ffast-math} and turn off
  floating-point contraction.
\end{itemize}

\emph{The schedule and the tilts.} The value $d_p$ of the schedule is evaluated once per prime in round-to-nearest
(with \code{libm}'s logarithm) and stored. The tilt used everywhere afterwards, in the random numbers $s_{p'}$ for
$p'>p$ and in $T_X$, is the double $\hat\nu_p=\mathrm{RU}\bigl(1/\mathrm{RD}(1-d_p)\bigr)\ge1/(1-d_p)$. We therefore
apply \cref{thm:criterion} and \cref{prop:bbmst61} with the schedule $\delta_p:=1-1/\hat\nu_p\ge d_p$, for which
$\nu_p=\hat\nu_p$ exactly. The bounds at level $p$ use $t\le d_p\le\delta_p$ in (i) and $d=\min(d_p,\frac12)\le\delta_p$
in (ii) and (iii), so they are admissible. All $\delta_p$ are at most $0.46+10^{-15}<\frac12$. The programs assert
$\hat\nu_p\le p$ at every prime, and $t\le d_p$ holds by construction (\code{rigx.c} uses $t=d_p$) or is asserted
(\code{rigcert.c}).

\emph{A rounding gap and its repair.} An audit found one rounding-direction flaw in \code{rigcert.c}: $(p-1)^2$ was
formed by converting an exact integer to a double, which can round \emph{up} in round-upward mode for
$p\ge94{,}906{,}267$, where $(p-1)^2>2^{53}$, although it is used as a divisor of upper bounds (in the second moments
and in the factors of $T_X$). These bounds could therefore be low by one unit in the last place. The program now uses
$\mathrm{RD}((p-1)^2)$ and reproduces the certified output for \cref{thm:lean} bit for bit
(\code{src/CERT\_NOTES.md}); \code{rigx.c} had the repair from the start.

\subsection{The certificate for \cref{thm:main}}\label{sec:certified}

The certified run for $m=\Mzero$ (command in \cref{app:reproduce}) prints
\[
  \eta_{\mathrm{up}}=\etaMainTwelve\qquad(A=\blockAMain,\ B=\blockBMain,\ C=\blockCMain),
\]
where $A$, $B$ and $C$ are the contributions of the primes $p<50$, $50\le p\le\PX$ and $\PX<p\le\PmaxMain$. The total
is the double $\etaMainFull$ (\code{\etaMainHex}), rounded up in the twelfth decimal. The program also prints
$k=\pi(\PmaxMain)=\kPrimesMain$, the upper bound $T_{\mathrm{up}}=\TupMain$ for $T_{\PmaxMain}$, and the quantities of the
test of \cref{prop:bbmst61}. \Cref{tab:budget} shows the budget by prime ranges, next to that of the certificate for
\cref{thm:lean}, and \cref{app:budget} gives finer breakdowns. About $\shareBelowTwoK\%$ of the budget is spent at the
primes below $2{,}000$.

\begin{table}[ht]
\centering
\caption{Where the budget goes. Certified contributions of the prime ranges for $m=\Mzero$
($P_{\max}=\PmaxMain$) and $m=\MzeroLean$ ($P_{\max}=\Pmax$); the second column gives the number of primes (for the
last two rows, for $m=\Mzero$ and for $m=\MzeroLean$). Every entry is rounded up in the sixth decimal; the totals
are the certified values of $\eta_{\mathrm{up}}$, rounded up.}
\label{tab:budget}
\small
% generated by make_figures.py from certificate-14501/output and logs/ -- do not edit
\begin{tabular}{@{}lrrr@{}}
\toprule
primes & number & $m=14{,}501$ & $m=16{,}000$ \\
\midrule
$p<17$ & 6 & 0.037108 & 0.034611 \\
$17\le p<50$ & 9 & 0.197793 & 0.195961 \\
$50\le p<200$ & 31 & 0.364896 & 0.365330 \\
$200\le p<2{,}000$ & 257 & 0.314659 & 0.312354 \\
$2{,}000\le p<20{,}000$ & 1{,}959 & 0.072697 & 0.071735 \\
$20{,}000\le p\le 2\cdot 10^5$ & 15{,}722 & 0.011319 & 0.010886 \\
$2\cdot 10^5<p\le P_{\max}$ & 50{,}829{,}550 / 11{,}060{,}953 & 0.001458 & 0.006511 \\
\midrule
total, $p\le P_{\max}$ & 50{,}847{,}534 / 11{,}078{,}937 & 0.999927 & 0.997385 \\
\bottomrule
\end{tabular}

\end{table}

\begin{proof}[Proof of \cref{thm:main}]
Apply \cref{prop:bbmst61} with $m=\Mzero$, $X=\PmaxMain$, the schedule $\delta_p=1-1/\hat\nu_p$ of \cref{sec:rounding}
for $p\le X$ and $\delta_p=\frac12$ for $p>X$, and $c_p$ the per-prime costs computed by the certified program. Each
$c_p$ is an upper bound for one of the quantities (i)--(iii) of \cref{thm:criterion}: for $p<17$ the first moment,
which is (i) with $t=0=\delta_p$, and otherwise a comparison bound with $t=d_p\le\delta_p$, multiplied by a margin
factor larger than $1$. The program certifies $\eta=\sum_{p\le X}c_p\le\etaMainTwelve<1$, and with
$\mu_{\mathrm{lo}}=\mathrm{RD}(1-\eta_{\mathrm{up}})\ge\muLoMain$ it certifies
\[
  \frac{T_X}{1-\eta}\le f_{\mathrm{up}}=\mathrm{RU}\Bigl(\frac{T_{\mathrm{up}}}{\mu_{\mathrm{lo}}}\Bigr)\le\fUpMain,
  \qquad
  (\log k+\log\log k-3)^2k\ge\critLoMain,
\]
where $k=\kPrimesMain\ge10$ and the logarithms are bounded below as in \cref{app:rounding}. Since
$\fUpMain<\critLoMain$, \cref{prop:bbmst61} applies. Therefore no family of congruences with distinct moduli, all at
least $\Mzero$, covers $\Z$. \Cref{cor:main} follows, since a covering system with distinct moduli must then have a
modulus below $\Mzero$.
\end{proof}

The test of \cref{prop:bbmst61} passes with $f/(\log k+\log\log k-3)^2k\approx0.611$; it would hold for any
$\eta$ up to about $0.999955$ (\code{certificate-14501/xcheck/results/tail.out}), so the margin of the certificate is
about $2.8\cdot10^{-5}$ in $\eta$. The margin of \cref{thm:criterion} is not enough here, as noted in \cref{sec:tail}.

\subsection{The formal certificate for \cref{thm:main}}\label{sec:certified-formal}

The formal proof of \cref{thm:main} (\cref{sec:lean}) uses a second checker, written in Lean (\code{Checker2Impl}). It
evaluates the per-prime bounds as in \cref{sec:SBH,sec:tausplit}, for $m=\Mzero$, the primes $p\le X=\XFormal$ and
$t=\delta_p$, in fixed-point arithmetic with values $n/2^{62}$ and explicit directed rounding. The schedule is that of
the C certificate with the two knots at $p\ge5\cdot10^4$ set to $0.45$, rounded to multiples of $10^{-9}$, with
$\delta_p=0$ for $p<17$ and $\delta_p=\frac12$ for $p>X$.
\begin{itemize}
\item For $p<17$ it bounds $\E[U_p(s)]$, which is (i) with $t=\delta_p=0$.
\item For $17\le p<m/2$ it uses \cref{lem:SBH} in the lower-tail form
  $\E[(\tau_ST-\beta)^+]=\tau_S\E[T]-\E[\beta]+\E[(\beta-\tau_ST)^+]$. The last term involves the joint law of $(T,b)$
  only for $T\le\beta/\tau_S$, so this evaluation has neither lost mass nor exponent truncation. The enumeration of
  $s_S$ is pruned as in \cref{sec:SBH}, at the relative level $2^{-33}$, and the pruning bounds are added.
\item For $m/2\le p\le\PX$ it uses the $\tau$-split (\cref{lem:tausplit}) with the truncations of \cref{sec:lost},
  with the law of $\tau_S$ kept on $[0,2^{22}]$.
\item For $\PX<p\le X$ it takes the primes in blocks of relative width $2^{-13}$ and evaluates the $\tau$-split
  bound once per block, at the first element of the block, where it is largest.
\end{itemize}
The run certifies
\[
  \sum_{p\le X}c_p\le\etaFormal,\qquad T_X\le\TupFormal;
\]
the primes $p<17$, $17\le p\le\PX$ and $\PX<p\le X$ contribute $\etaFormalA$, $\etaFormalB$ and $\etaFormalC$.

\begin{proof}[Formal proof of \cref{thm:main}]
Apply \cref{thm:criterion} in the form of \cref{rem:criterion-sharp}, with $m=\Mzero$, $X=\XFormal$, the schedule above,
the per-prime costs $c_p$ of the Lean checker (each an upper bound for the quantity (i) with $t=\delta_p$), and the
constant $C_X=\frac{5941}{50000}$ of \cref{prop:tailtwo}. The left-hand side of the criterion is at most
\begin{multline*}
  \etaFormal+\frac{5941}{50000}\cdot\frac{\TupFormal}{\XFormal}\\
  \le\etaFormal+\tailFormal\le\totalFormal<1 .
\end{multline*}
Therefore no family of congruences with distinct moduli, all at least $\Mzero$, covers $\Z$.
\end{proof}

The whole argument, including the soundness of every step of the checker and the tail estimate, has been checked by
Lean's kernel. The only trusted step is the compiled evaluation of the checker (\cref{sec:lean}), which takes about ten
minutes and $2.6$ GB. The margin is $\slackFormal$. For $m=14{,}500$ a floating-point prototype of the same checker
gives $\sum_pc_p>1$ on its own, so $\Mzero$ is also the limit of the formal version of the method. On the way, the
same development certified $m=14{,}600$ with $X=\Pmax$ and the tail of \cref{prop:tail} (total $\totalLeanTwo$), and
$m=14{,}505$ with $X=10^9$ (total $\totalLeanThree$). The C program \path{certificate-14501/src/rigx.c}, run at
$m=14{,}600$ and $X=\Pmax$ with matching settings, agrees with the Lean checker's sum $\sum_{p\le X}c_p$ to within
$10^{-6}$.

\subsection{The certificate for \cref{thm:lean}}\label{sec:certified-lean}

The certified run for $m=\MzeroLean$ prints
\[
  \eta_{\mathrm{up}}=\etaCertTwelve\qquad(A=\blockA,\ B=\blockB,\ C=\blockC),
\]
for the primes $p<50$, $50\le p\le\PX$ and $\PX<p\le\Pmax$. The total is the double $\etaCertFull$
(\code{\etaCertHex}), rounded up in the twelfth decimal. The program also prints $k=\pi(\Pmax)=\kPrimes$ and the upper
bound $T_{\mathrm{up}}=\Tup$ for $T_{\Pmax}$.

\begin{proof}[Proof of \cref{thm:lean}]
Apply \cref{thm:criterion} with $m=\MzeroLean$, $X=\Pmax\ge2^{27}$, the schedule $\delta_p=1-1/\hat\nu_p$ of
\cref{sec:rounding} for $p\le X$ and $\delta_p=\frac12$ for $p>X$, and $c_p$ the per-prime costs computed by the
certified program, each of which is an upper bound for one of the quantities (i)--(iii). The program certifies
$\sum_{p\le X}c_p\le\etaCertTwelve$ and $T_X\le T_{\mathrm{up}}<\TupCeil$. Hence the left-hand side of
\eqref{eq:criterion} is at most
\[
  \etaCertTwelve+\frac{100}{189}\cdot\frac{\TupCeil}{\Pmax}
  \le\etaCertTwelve+\tailElem<\etaElem<1 .
\]
Therefore no family of congruences with distinct moduli, all at least $\MzeroLean$, covers $\Z$.
\end{proof}

The criterion of BBMST gives an independent confirmation of the tail. With
$\mu_{\mathrm{lo}}=\mathrm{RD}(1-\eta_{\mathrm{up}})\approx\muLo$, the program certifies
$T_X/(1-\eta)\le f_{\mathrm{up}}\le\fUp$ and $(\log k+\log\log k-3)^2k\ge\critLo$, so \cref{prop:bbmst61} applies as
well, with $f/(\log k+\log\log k-3)^2k\approx0.043$ (\code{logs/xcheck\_final.md}).

\subsection{Independent recomputation}\label{sec:xcheck}

\emph{The certificate for \cref{thm:main}.} Two independent checks were made.
\begin{enumerate}[label=(\arabic*)]
\item \emph{Cross-check by independently written code} (\code{certificate-14501/xcheck/}). The programs of the
  recomputation of the certificate for \cref{thm:lean} described below, which were written independently of both
  checkers, were rerun with the new schedule (only their schedule header was changed), together with a new program
  for the large primes that uses no code of \code{rigx.c}. For each prime they compute a rigorous interval
  $[\mathrm{lo},\mathrm{hi}]$ for the exact value of the level bound. \Cref{tab:xcheck-main} summarizes the results:
  the certified cost is at least $\mathrm{hi}$ at every prime compared, and the total of the upper ends for
  $\PX<p\le\PmaxMain$ is below the certified $C$.
\item \emph{Independent referee.} An AI agent acting as a referee checked the mathematics of
  \cref{sec:SBH,sec:tausplit,sec:lost} by hand; brute-forced the identity of \cref{lem:SBH} with exact rationals on
  $36{,}840$ random pairs $(p,s)$ at all $\nPrimesSBH$ primes where it is used, including cases with $y$ itself a
  prime of $B$ and with $v_q\ge2$ for $q\in B$ (no mismatch); audited the rounding directions and the lost-mass
  accounting; and recomputed every level bound for all primes up to $\PmaxMain$ with separately written programs (an
  exact complement identity with a pruned enumeration for $p<\firstTau$, and a logarithmic-grid law of $\tau(s)$ with
  its own sieve for the larger primes), with exact-rational and $60$-digit spot checks. It obtained the rigorous
  interval $[\refEtaLo,\refEtaHi]$ for the sum of the level bounds, found the certified cost at or above its upper
  bound at every prime (at every level up to $\PX$ and at every $500$-th prime beyond, where the per-prime costs are
  recorded), reproduced the run bit for bit on macOS with Apple clang and with GNU gcc, and found no errors. Its remaining
  remarks were that the trust base includes \cite[Theorem~6.1]{BBMST} and that the bound for the largest primes is
  conservative (\cref{sec:tausplit}).
\end{enumerate}

\begin{table}[ht]
\centering
\caption{The cross-check of the certificate for \cref{thm:main}. Here $\mathrm{hi}$ is the upper end of an
independent rigorous interval for the exact level bound, and ``cost'' is the certified cost, including the margins of
\cref{sec:eval-overview}.}
\label{tab:xcheck-main}
\small
\begin{tabular}{@{}>{\raggedright\arraybackslash}p{0.3\textwidth}>{\raggedright\arraybackslash}p{0.64\textwidth}@{}}
\toprule
check & result\\
\midrule
schedule, all $\nPrimesPX$ primes $p\le\PX$ & relative deviation from a $40$-digit evaluation at most $4.5\cdot10^{-16}$;
  $\hat\nu_p\ge1/(1-d_p)$ everywhere\\
first moments, $p<17$ & exact rational values; at most the certified cost at all $6$ primes\\
S/B/H levels, $17\le p\le\lastSBH$ & all $\nPrimesSBH$ primes, with two different splits whose intervals always overlap;
  $0$ violations; cost/hi $\in[1+9.2\cdot10^{-7},\,1+9.6\cdot10^{-7}]$\\
$\tau$-split levels, $\firstTau\le p\le\PX$ & all $\nPrimesTauPX$ primes, with a direct law of $\tau(s)$; $0$ violations;
  cost/hi $\in[1+4.5\cdot10^{-5},\,1+1.2\cdot10^{-4}]$\\
$\tau$-split levels, $\PX<p\le\PmaxMain$ & $118{,}717$ primes compared (all $p\le\PX$ and every $500$-th prime beyond),
  $0$ violations; the upper ends summed over all $\nPrimesCMain$ primes give $0.0014578122\le C$\\
tail & $k=\kPrimesMain$ confirmed by a segmented sieve; $T_k=705974.98819$, below $T_{\mathrm{up}}$ by $1.0\cdot10^{-8}$
  relative; $f/(\log k+\log\log k-3)^2k=0.611$\\
\bottomrule
\end{tabular}
\end{table}

\emph{The certificate for \cref{thm:lean}.} Separately written code (\path{tests/xcheck_final/}), with its own
decomposition and its own outward rounding, recomputed this certificate (\path{logs/xcheck_final.md}); it reuses none
of the checker's machinery. For each of the $\nPrimesPX$ primes $p\le\PX$ and each candidate bound of the checker, it
computed a rigorous interval $[\mathrm{lo},\mathrm{hi}]$ for the exact value and checked
$c_p\ge\min_{\text{candidates}}\mathrm{hi}$ in exact rational arithmetic: there were $0$ violations and $0$
inconclusive cases, and the smallest relative margin, $1.4\cdot10^{-7}$ at $p=8009$, is seven times the width of the
interval there. For $31\le p<8000$ it uses the identity of \cref{lem:SBH}, and for $8000<p\le\PX$ a divisor-count
dynamic program on $[1,2^{20}]$ and the exact identity $\E[(\tau-x)^+]=\E\tau-x+\sum_{T\le x}(x-T)\P(\tau=T)$, which
needs no truncation; both are run in both rounding directions and agree with brute force on small instances. It also
reproduced the schedule bit for bit, the first and second moments (against $60$-digit values and direct double sums),
the total $0.0065107856047$ of the fractional-moment bounds over $(\PX,\Pmax]$ (against the checker's $\blockC$), the
prime count, the product $T_k=321051.3638366$ (the checker's $T_{\mathrm{up}}$ is larger by $2.2\cdot10^{-9}$ relative)
and the tail criterion. It also measures the slack of the cruder evaluation. The sum over $p\le\PX$ of the exact values of the
best candidate bounds is $\xcheckExact$, against the checker's $\sumCostPX$; the difference, about $\xcheckSlack$, is
almost entirely at $p\ge50$. The recomputations check numbers, not mathematics: the comparison theorem and the
reduction of \cref{sec:loss} are outside their scope. Their interval code relies on the compiler honouring directed
rounding (it passes its own self-tests and was cross-validated against brute force and $60$-digit arithmetic).

\subsection{Tests}\label{sec:tests}

\emph{The checkers' formulas.} On small instances, independent Python programs compare every quantity that
\code{rigcert.c} computes with brute-force enumeration of $s$ (with a rigorous bracket for the pruned part), with exact
rational arithmetic (for $\E[U_p]$ and $\E[U_p^2]$, via a different decomposition) and with $50$-digit arithmetic (for
the products $\E[\tau(s)^\theta]$) (\code{tests/cert/run\_small.sh}). The instances use the knot schedule and the
constant band of the final run, values $\delta>\frac12$, and tiny truncation parameters that force every tail path.
There were no failures over $308$ comparison-bound candidates, $84$ moment-bound candidates and the first and second
moments at every prime; on two larger instances, $108$ primes of exact moments and $828$ pairs $(p,\theta)$ of
$50$-digit products agreed. Earlier validations on other configurations, including $431$ exact prime checks, are
described in \code{src/CERT\_NOTES.md}. For \code{rigx.c}, plain enumeration of $s$ was compared with $42$ cases at
$m=300$ and $m=1000$, including settings that force every truncation and lost-mass path
(\code{certificate-14501/validation/}); the certified values were always at least the lower end of the brute-force
bracket, and inside the bracket with the default settings. At $m=\MzeroLean$ with the schedule of \cref{thm:lean},
\code{rigx.c} agrees with the exact intervals of the recomputation above to within $10^{-9}$ relative at the S/B/H
levels and lies inside the intervals at the others.

\emph{The comparison theorem.} About $10^6$ randomized, exhaustive and adversarial tests (\code{tests/review/})
found no violation of \cref{lem:onecoord}, \cref{thm:comparison} and \cref{lem:pointwise}. They compare the
theorem with the exact worst case over all $\nu$-bounded product measures (computed by a backward bathtub recursion),
include an exhaustive search over all families of at most three cylinders on $\Z/4\times\Z/3$, and simulate the
measures of BBMST exactly on small moduli, checking \cref{lem:F1}, the kernel bounds of \cref{lem:kernel},
\cref{lem:F3} and each inequality in the proof of \cref{thm:perprime}.

\emph{Proof review.} Separate AI referee passes checked each step of the proof of the comparison theorem and of the
reduction, and the new steps of \cref{sec:SBH,sec:tausplit,sec:lost}; they found no gap, and their presentational
suggestions have been incorporated. No human has refereed the argument yet.

\subsection{Run times and reproducibility}\label{sec:runtimes}

The certified run for \cref{thm:main} takes $87$--$93$ seconds on one core of an Apple M3 Max and about $1.2$~GB of
memory (for the sieve up to $10^9$). Rebuilt from the published source, it reproduces the certified result lines and
the per-prime dump bit for bit; the referee reproduced them on macOS with Apple clang and with GNU gcc. On Linux
(x86-64, glibc, GNU gcc at \code{-O0} and at \code{-O2}, which agree) the run also certifies, with the same result
line; the last digits of the certified value and $503$ of the $119{,}644$ per-prime costs in the dump differ, by at
most about $2\cdot10^{-12}$ relative, because of the logarithm (see below). The cross-check of
\cref{tab:xcheck-main} takes about $29$ minutes on one core. The certified run for \cref{thm:lean} takes $156.6$
seconds of CPU time on one core (\code{logs/timing.txt}) and is also reproduced bit for bit
(\code{logs/repro16k.txt}). Bit-for-bit reproduction requires floating-point contraction to be off and the same
\code{libm} logarithm, which is used only to evaluate the schedule; another \code{libm} may move some $d_p$ by one
unit in the last place. The last digits of the output would then change, but the certificates would remain valid,
since any schedule is admissible and the checkers use their own values consistently.

\subsection{Formal verification}\label{sec:lean}
% Status (2026-09-28): complete for thm:main, thm:lean, thm:squarefree, thm:sf23, thm:sf7type, thm:odd, thm:odd2;
% source DELIV/lean/README.md (coordinator's axiom audits).

\Cref{thm:main,thm:lean} have been formally verified in Lean~4 (v4.34.1) with Mathlib, each modulo one
\code{native\_decide} evaluation. The development is in the folder \path{lean/} of the accompanying files, and
\path{lean/README.md} describes the build and the audit. The root theorems, in
\path{lean/MinModulus/Checker2Sound/Final14501.lean} and \path{lean/MinModulus/Main/Main.lean}, are
\begin{alltt}\small
theorem MinModulus.Checker2Sound.not_covers_14501 \{\(\iota\) : Type*\} [Fintype \(\iota\)]
    (d : \(\iota\) \(\to\) \(\mathbb{N}\)) (a : \(\iota\) \(\to\) \(\mathbb{Z}\)) (hd : Function.Injective d)
    (hm : \(\forall\) i, \MzeroNum \(\le\) d i) :
    \(\exists\) x : \(\mathbb{Z}\), \(\forall\) i, \(\neg\) ((d i : \(\mathbb{Z}\)) \(\mid\) x - a i)
theorem MinModulus.not_covers \{\(\iota\) : Type*\} [Fintype \(\iota\)]
    (d : \(\iota\) \(\to\) \(\mathbb{N}\)) (a : \(\iota\) \(\to\) \(\mathbb{Z}\)) (hd : Function.Injective d)
    (hm : \(\forall\) i, \MzeroLeanNum \(\le\) d i) :
    \(\exists\) x : \(\mathbb{Z}\), \(\forall\) i, \(\neg\) ((d i : \(\mathbb{Z}\)) \(\mid\) x - a i)
\end{alltt}
Neither contains \code{sorry}. For each, \code{\#print axioms} lists only \code{propext}, \code{Classical.choice},
\code{Quot.sound} and one auxiliary axiom introduced by \code{native\_decide}:
\begin{center}\small\code{MinModulus.Tail2.check2T14501\_eq\_true.\_native.native\_decide.ax\_1\_1}\\
\code{MinModulus.CheckerImpl.check\_eq\_true.\_native.native\_decide.ax\_1\_1}\end{center}
respectively. Each states that one closed Boolean (\code{check2T14501}, respectively \code{check}) evaluates to
\code{true}. The trust base is therefore Lean's kernel together with the compiled evaluation of this Boolean. That
takes about ten minutes and $2.6$ GB for \code{check2T14501}, and about $50$ seconds for \code{check}. It trusts
Lean's compiler and runtime, including the arbitrary-precision arithmetic on natural numbers. The code of both
checkers contains no \code{implemented\_by}, \code{extern}, \code{unsafe} or \code{partial} declarations; the only
compiler rewrites are two \code{@[csimp]} lemmas, which are proved equalities. The evaluation of \code{check} has been
re-run from a fresh file. The whole development, including every native evaluation, has also been rebuilt from a fresh
copy of the sources (with the Mathlib build shared), in about eleven minutes. Both gave the same results and the
same axioms.

\emph{The formalization of \cref{thm:main}.} It reuses the modules \code{Arith}, \code{Distortion},
\code{Comparison}, \code{Smooth}, \code{Tail} and \code{Main} of the first formalization unchanged, including the
certificate interface \code{Cert}. It also reuses the checker's arithmetic, sieve and first-moment code with their
soundness proofs. New are the following parts.
\begin{itemize}
\item The checker \code{Checker2Impl}, which implements \cref{sec:certified-formal}.
\item The mathematics of its quantities, in \code{Checker2Math}:
  \begin{itemize}
  \item \cref{lem:SBH,lem:tausplit} and the domination $\tau(s)\le2^{\sum v_q}$;
  \item the lower-tail identity;
  \item the joint law of $(T,b)$ and the law of the exponent sum;
  \item the depth-first enumeration and its pruning bound.
  \end{itemize}
\item A code-level soundness proof, \code{Checker2Sound}, covering:
  \begin{itemize}
  \item the dynamic programs, their states and the tables;
  \item the S/B/H and $\tau$-split costs;
  \item the loops and the blocks;
  \item the assembly of the certificate interface (\code{E.run2\_sound}).
  \end{itemize}
\item The sharpened tail estimate \cref{prop:tailtwo} and the corresponding form of \cref{thm:criterion}
  (\code{Tail2}), and the certificate \code{check2T14501\_eq\_true}.
\end{itemize}
The soundness theorem \code{E.run2\_sound} holds for all parameters of the checker that satisfy four decidable side
conditions. A certificate for another $m$, or another schedule, therefore needs only a new \code{native\_decide}
evaluation.

\emph{Structure of the first formalization.} The whole development has $255$ files and about $52{,}700$ lines,
counting the formalizations of \cref{thm:main,thm:lean,thm:squarefree,thm:sf23,thm:sf7type,thm:odd,thm:odd2} but not
the $194$ generated modules that embed and check the certificate data of \cref{thm:sf7type}. Of these, $64$ files and
about $18{,}400$ lines belong to the first formalization. \Cref{tab:lean} lists the modules and their correspondence with this paper. The first formalization
follows the paper, with two differences: it uses the
comparison theorem only for the hinge functions, since the other bounds are obtained from the hinge with
$t=\delta_p$ by pointwise inequalities on the aligned side, and it treats the primes above $X=\Pmax$ by
\cref{prop:tail}, so that it needs neither \cite[Theorem~6.1]{BBMST} nor Dusart's estimate.

\emph{The checker of the first formalization.} The numerical certificate for \cref{thm:lean} is produced by an
independent implementation of the checker in core Lean (\code{CheckerImpl}), not a transcription of
\code{src/rigcert.c}. It works in fixed point, with values $n/2^{62}$ and explicit directed rounding; it uses
$t=\delta_p$ only, with $\delta_p$ an exact rational (a multiple
of $10^{-9}$ obtained from the same knots) and exact rational tilts; and for $\PX<p\le\Pmax$ it processes the primes
in blocks, taking the factors of each block at its first prime, with $\theta\in\{2,\frac52,3\}$. It certifies
$\sum_{p\le\Pmax}c_p\le\etaLean$ and $T\le\TupLean$ (slightly larger than in the C program, because of the blocks),
so that the total with the tail term $\frac{100}{189}T/X=\tailLean$ is $\totalLean<1$. A code-level soundness proof
(\code{CheckerSound}) shows that the checker's loops and arrays bound the true per-prime losses (invariants of the
dynamic programs, semantics of the enumeration, the blocks) and assembles them into the certificate
\code{Cert \MzeroLeanNum{} (2*10\^{}8) \(\delta\) c T} consumed by \code{MinModulus.Main.not\_covers\_of\_cert}, whose
conditions correspond to those of \cref{thm:criterion} with $t=\delta_p$. The squarefree theorem of
\cref{sec:squarefree} is formalized in the same way, with its own checker. \Cref{thm:sf23} is formalized as
\code{MinModulus.SliceCert.not\_covers\_sf3\_23smooth}, an instance of a general theorem
\code{not\_covers\_of\_slices} for certificates of the kind described in \cref{sec:squarefree} (any start moduli and
residual primes), which uses only the three standard axioms; the certificate itself is checked by one
\code{native\_decide} evaluation (\code{check23\_eq\_true}, about $70$ seconds).

\emph{The formalization of \cref{thm:sf7type}.} It is the theorem
\code{MinModulus.SquarefreeM7.not\_covers} in the module \code{SquarefreeM7}; its hypotheses are those of
\code{MinModulus.not\_covers} above, together with \code{Squarefree (d i)} and $7\le{}$\code{d i}. It has four parts, which interact only through propositions fixed in advance.
\begin{itemize}
\item The reduction to the nine types, the symmetries and the classes modulo $30$ (\code{Reduction}), checked
  entirely by the kernel.
\item The per-level bound for families with repetitions at every level, with the tail of \cref{thm:sf7type}: the
  levels $7$, $11$ and $13$ with the Jensen splits and the exact terms (including the split by the class of the
  modulus $390$), the grouped bound with power-of-two thresholds for the other levels up to $2^{31}$, and
  \cref{prop:tail} beyond (\code{Bound}).
\item The tail coefficients: a sieve up to $2^{31}$ and the Poisson--binomial recursions with upward rounding, with a
  proof that they bound the true sums, run in $16$ chunks for each of the two schedules (\code{Coef}).
\item A checker in core Lean with exact integer arithmetic, about $1.6$ times slower than the C verifier, with its
  soundness proof, and the nine certificates, converted to an exact format and embedded as data ($446$~MB in $83$
  chunks; \code{Check}, \code{Data}, \code{Cert}).
\end{itemize}
\code{\#print axioms} lists \code{propext}, \code{Classical.choice}, \code{Quot.sound} and $124$ auxiliary axioms of
\code{native\_decide}: $83$ for the certificate chunks, $9$ for the level-$7$ covers of the nine types and $32$ for the
tail coefficients. The reduction and the analytic part use only the three standard axioms. A build of the whole
development from scratch takes about $80$ minutes on $24$ threads, almost all of it in the certificate chunks.

\emph{What is not formalized.} The C certificate of \cref{thm:main} (\cref{sec:certified}) and
\cite[Theorem~6.1]{BBMST}, which it uses, are not formalized; the formal proof does not need them. \Cref{thm:odd} is
formalized (\cref{sec:odd}), and so is \cref{thm:odd2}; \cref{thm:sf29,thm:sf6smooth} and the results of
\cref{sec:rings} are not.

\begin{table}[ht]
\centering
\caption{Modules of the Lean formalization (\code{lean/MinModulus/}) and the corresponding results of this paper.}
\label{tab:lean}
\small
\begin{tabular}{@{}lr>{\raggedright\arraybackslash}p{0.69\textwidth}@{}}
\toprule
module & lines & contents\\
\midrule
\code{Arith} & $679$ & setup of a family of congruences: $Q$, the primes and levels, the Chinese remainder
  theorem, the sets $B_\ell$, \cref{lem:unionbound,lem:pointwise,lem:enlargement}\\
\code{Distortion} & $646$ & the measures \eqref{eq:rule}, kernels and tilts, \cref{lem:F1,lem:F2,lem:F3},
  \cref{prop:F4,lem:kernel}\\
\code{Comparison} & $710$ & \cref{lem:greedy,lem:onecoord} (by duality), \cref{prop:rectangles,lem:vlaw}, and
  \cref{thm:comparison} for hinges\\
\code{Smooth} & $2142$ & the tilted law of $s$, $U_p(s)$, the moments $\E[\tau(s)^\theta]$, hinge bounds, monotone
  coupling in the exponent caps\\
\code{Tail} & $354$ & \cref{prop:tail}\\
\code{Main} & $613$ & the certificate interface \code{Cert}, \code{not\_covers\_of\_cert} (\cref{thm:perprime} with
  $t=\delta_p$, \cref{thm:criterion}), the final theorem\\
\code{CheckerImpl} & $1269$ & the executable checker and \code{check\_eq\_true} (by \code{native\_decide})\\
\code{CheckerMath} & $4094$ & the mathematics of the checker's quantities: laws of $\tau(s_A)$ and $\tau(s_B)$,
  the split of \cref{lem:AB}, the comparison bound, first and fractional moments\\
\code{CheckerSound} & $6201$ & code-level soundness: loop and array invariants, enumeration, dynamic programs,
  blocks, and the assembly of the certificate (\code{CheckerSound.D.cert\_16000})\\
\code{Squarefree} & $1649$ & \cref{thm:squarefree}, with its own checker and soundness proof\\
\midrule
\multicolumn{3}{@{}l}{\emph{formalization of \cref{thm:main}}}\\
\code{Checker2Impl} & $1139$ & the executable checker of \cref{sec:certified-formal}\\
\code{Checker2Math} & $2155$ & \cref{lem:SBH,lem:tausplit} in the form of the hinge loss, the lower-tail identity, the
  joint law of $(T,b)$, the law of the exponent sum, the enumeration and its pruning bound\\
\code{Checker2Sound} & $6604$ & code-level soundness: dynamic programs, tables, the S/B/H and $\tau$-split costs,
  loops and blocks, the generic soundness theorem \code{E.run2\_sound}, and the final theorem
  \code{not\_covers\_14501}\\
\code{Tail2} & $2022$ & \cref{prop:tailtwo}, the criterion with the sharpened tail, and the certificate
  \code{check2T14501\_eq\_true} (by \code{native\_decide})\\
\midrule
\code{Odd} & $4381$ & \cref{thm:odd,thm:odd2}: the chain for restricted moduli, the per-case weights, a checker for
  the $\tau$-split bounds with excluded and capped primes, its soundness, and nine \code{native\_decide}
  certificates\\
\code{SliceCert} & $3138$ & \cref{thm:sf23}: the class decomposition modulo $30$, the per-level bound for families
  with repetitions, certificate records and the branch-and-bound tree with symmetries, the start-case check, and
  \code{check23\_eq\_true} (by \code{native\_decide})\\
\code{SquarefreeM7} & $14950$ & \cref{thm:sf7type}: the reduction to the nine types, the per-level bound with the
  refinements at $7$, $11$ and $13$ and the tail, the tail coefficients up to $2^{31}$, the exact checker and its
  soundness (and $194$ generated modules: the certificate data and $124$ \code{native\_decide} checks)\\
\bottomrule
\end{tabular}
\end{table}
% CHECK: the Lean checker's totals are from DELIV/lean/README.md and SCRATCH/agents/checkerimpl/full_final.out
%        (0.997384786285, 321095.49941797, tail 0.000849458994, total 0.998234245279); they are not under DELIV/logs.
